There are three common ways to traverse the nodes of a binary tree:

\textbullet\ Preorder: First process the root, then the left subtree, and finally the right subtree.
\textbullet\ Inorder: First process the left subtree, then the root, and finally the right subtree.
\textbullet\ Postorder: First process the left subtree, then the right subtree, and finally the root.

There is a binary tree of $n$ nodes with distinct labels. You are given the preorder and inorder traversals of the tree, and your task is to determine its postorder traversal.

\vspace{0.3cm}\noindent\textbf{Input}

The first input line has an integer $n$: the number of nodes. The nodes are numbered $1,2,\dots,n$.
After this, there are two lines describing the preorder and inorder traversals of the tree. Both lines consist of $n$ integers.
You can assume that the input corresponds to a binary tree.

\vspace{0.3cm}\noindent\textbf{Output}

Print the postorder traversal of the tree.

\vspace{0.3cm}\noindent\textbf{Constraints}

\textbullet\ $1 \le n \le 10^5$